\documentclass{scrbook}
\usepackage{geometry}
\geometry{paperwidth=8in, paperheight=10in, inner=0.875in, outer=1.75in, top=0.875in, bottom=0.875in, marginparwidth=1.25in, marginparsep=0.25in}
\usepackage{fontspec}
\setmainfont{Palatino}
\setsansfont{Helvetica Neue}[Scale=0.92]
\usepackage{xcolor}
\usepackage{tikz}
\usepackage{tikzpagenodes}
\usepackage{hyperref}
\usepackage{lettrine}
\definecolor{accentcolor}{HTML}{1A4D3E}

\usepackage[explicit]{titlesec}
\newcommand{\mlsysChapNumFont}{\normalfont\rmfamily\fontsize{70pt}{70pt}\selectfont\bfseries}

\titleformat{\chapter}[block]
  {\normalfont\rmfamily\bfseries\color{accentcolor}}
  {}
  {0pt}
  {%
    \begin{tikzpicture}[remember picture, overlay]
      \node[anchor=north east,
            inner sep=0pt,
            font={\mlsysChapNumFont},
            text=accentcolor]
        at ([xshift=-40pt]current page.north east |- current page text area.north)
        {\thechapter};
    \end{tikzpicture}%
    {\fontsize{34pt}{38pt}\selectfont #1}%
  }
  []

\newcommand{\chaptersubtitle}[1]{\par\vspace{0.4em}{\normalfont\normalsize\itshape #1}}

\begin{document}
\chapter{\texorpdfstring{The Causal Boundary\chaptersubtitle{Introduction to Physical AI Systems}}{The Causal Boundary}}\label{sec-boundary}

\renewcommand{\LettrineFontHook}{\rmfamily\bfseries\color{accentcolor}}
\renewcommand{\LettrineTextFont}{\normalfont\scshape}

\lettrine{E}{very} major epoch in computing has been defined by the boundary across which bits travel. In personal computing, bits crossed from mainframe backplanes to human eyes through glowing cathode-ray displays. In cloud computing, bits crossed global fiber-optic networks to synchronize distributed relational databases.

\end{document}
